\chapter{Proposed Solution and System Design}
\label{ch:pfe-design}

\section{Introduction}

The proposed solution is a static, evidence-gated pipeline that transforms a
compiled Android APK into a CycloneDX SBOM and a confidence-aware
vulnerability report. The design separates extraction, reference construction,
matching, version selection, and downstream security interpretation so that
each claim can be traced to its supporting evidence.

This separation is not only architectural; it is also epistemic. The system
must preserve the boundary between what is observed in the APK, what is
reconstructed from the reference database, and what is merely suggested for
reporting or investigation. If those layers are merged too early, the final
report becomes difficult to audit and easy to over-interpret. The solution
therefore uses explicit data stores, explicit decision stages, and explicit
report-only experiment paths so that the final SBOM stays within the evidence
available to it.

\section{Functional and Non-Functional Requirements}

The main functional requirements are APK inspection, component candidate
generation, exact or ambiguous version estimation, CycloneDX generation,
vulnerability lookup, and human-readable reporting. The main non-functional
requirements are explainability, reproducibility, conservative claims,
bounded experimental isolation, and traceable evidence.

These requirements guide the implementation choices that appear later in the
chapter. The extractor must tolerate malformed or partial APK content without
silently discarding evidence. The matcher must report confidence and
alternatives instead of only a single winner. The reporting layer must
present analyst-readable outputs without changing authoritative component or
version decisions. Finally, the experimental layer must remain separable from
the calibrated path so that exploratory weights, R8 diagnostics, and AI
assistance cannot override the production evidence policy.

\section{General Architecture}

\begin{figure}[H]
\centering
\resizebox{\textwidth}{!}{%
\begin{tikzpicture}[
  core/.style={
    rectangle,
    rounded corners,
    draw=schoolblue,
    thick,
    fill=blue!5,
    align=center,
    text width=31mm,
    minimum height=11mm,
    font=\small
  },
  store/.style={
    cylinder,
    shape border rotate=90,
    aspect=0.28,
    draw=schoolblue,
    thick,
    fill=cyan!7,
    align=center,
    text width=30mm,
    minimum height=13mm,
    font=\small
  },
  output/.style={
    core,
    draw=evidencegreen,
    fill=green!6
  },
  research/.style={
    rectangle,
    rounded corners,
    draw=warningorange,
    dashed,
    thick,
    fill=orange!6,
    align=center,
    text width=32mm,
    minimum height=11mm,
    font=\small
  },
  arrow/.style={-{Latex[length=2.2mm]},thick,draw=schoolblue},
  outarrow/.style={-{Latex[length=2.2mm]},thick,draw=evidencegreen},
  researcharrow/.style={-{Latex[length=2.2mm]},thick,dashed,
    draw=warningorange}
]
\node[core] (apk) at (0,4.8) {Compiled Android APK};
\node[core] (inspect) at (4.3,4.8) {Archive and manifest inspection\\DEX and native discovery};
\node[core] (features) at (8.6,4.8) {DEX evidence extraction\\classes, methods, opcodes, APIs, literals, graphs};

\node[core] (maven) at (0,2.1) {Maven JAR/AAR\\coordinates and versions};
\node[core] (reference) at (4.3,2.1) {Download, D8 conversion,\\reference feature extraction};
\node[store] (db) at (8.6,2.1) {Versioned SQLite\\reference-feature DB};

\node[core] (candidates) at (12.9,4.8) {Candidate generation and\\multi-feature scoring};
\node[core] (decision) at (12.9,2.1) {Identity and version decisions\\confidence, ambiguity, abstention};
\node[output] (sbom) at (17.2,4.8) {CycloneDX 1.5 SBOM\\PURLs, evidence, native hashes};
\node[output] (osv) at (17.2,2.1) {OSV vulnerability lookup\\certainty-aware findings};
\node[output] (report) at (17.2,0.0) {JSON, Markdown, and HTML\\evidence report};
\coordinate (authorityright) at (19.4,2.4);

\node[research] (r8) at (4.3,-2.5) {R8 and multi-scale diagnostics\\basic blocks, call chains, structural gates};
\node[research] (ml) at (8.6,-2.5) {ML/DL assistance\\Random Forest and Transformer experiments};
\node[research] (router) at (12.9,-2.5) {Hybrid router and version-band\\report-only recommendations};

\node[
  draw=schoolblue,
  rounded corners,
  inner sep=4mm,
  fit=(apk)(inspect)(features)(maven)(reference)(db)(candidates)(decision)
      (sbom)(osv)(report)(authorityright),
  label={[schoolblue,font=\bfseries]above:Authoritative APK-to-SBOM path}
] {};
\node[
  draw=warningorange,
  dashed,
  rounded corners,
  inner sep=4mm,
  fit=(r8)(ml)(router),
  label={[warningorange,font=\bfseries]below:Report-only research and diagnostic layer}
] {};

\draw[arrow] (apk) -- (inspect);
\draw[arrow] (inspect) -- (features);
\draw[arrow] (features) -- (candidates);
\draw[arrow] (maven) -- (reference);
\draw[arrow] (reference) -- (db);
\draw[arrow] (db) -- (decision);
\draw[arrow] (candidates) -- (decision);
\draw[outarrow] (decision) -- (sbom);
\draw[outarrow] (decision) -- (osv);
\draw[outarrow] (sbom.east) -- (authorityright |- sbom.east) |- (report.east);
\draw[outarrow] (osv) -- (report);

\draw[researcharrow] (decision.south) -- ++(0,-8mm) -| (r8.north);
\draw[researcharrow] (r8) -- (ml);
\draw[researcharrow] (ml) -- (router);
\draw[researcharrow] (router.north) -- ++(0,5mm) -| (report.south);
\end{tikzpicture}
}
\caption{Detailed architecture of the implemented APK-to-SBOM system,
including the authoritative analysis and reference-building paths and the
non-authoritative ML/DL/R8 assistance layer.}
\label{fig:pfe-high-level-architecture}
\end{figure}

The architecture has two authoritative input paths. The upper path extracts
archive, manifest, DEX, native, and forensic evidence from the target APK. The
middle path converts versioned Maven artifacts into the same DEX-oriented
feature representation and stores them in SQLite. These paths meet at
candidate scoring and the separate identity/version decision stage. Accepted
decisions feed CycloneDX, OSV, and human-readable reports. The dashed orange
layer contains R8 diagnostics, multi-scale structural features, Random Forest
and Transformer experiments, and the hybrid router. It may annotate a report,
but it cannot change selected versions, SBOM components, vulnerability labels,
or source-validation metrics.

\section{APK Extraction and DEX Analysis}

The extraction layer discovers standard and nonstandard DEX files, parses
application metadata, records resource and asset inventory, and inventories
native files. The DEX analysis layer derives classes, methods, opcode
statistics, API references, string literals, and dependency structures used
during matching.

\section{Reference Database Design}

Reference features are built from versioned Maven JAR and AAR artifacts. Each
artifact is downloaded, converted to a comparable DEX representation when
necessary, analyzed, and stored in SQLite. Shared calibrated databases are
kept separate from evaluation-only and experiment-only databases.

\section{Evidence Model}

The matcher combines several evidence families:

\begin{table}[H]
\centering
\caption{Evidence families used by the component-attribution pipeline.}
\label{tab:pfe-evidence-families}
\begin{tabularx}{\textwidth}{L{0.27\textwidth}X}
\toprule
\textbf{Evidence family} & \textbf{Design role}\\
\midrule
Packages, classes, methods & Candidate generation and identity support\\
Opcode features & Bytecode similarity under partial identifier changes\\
Class dependency graph & Structural evidence across related classes\\
Literals and unique features & Discrimination between nearby versions\\
API/member references & Optional contextual signal\\
POM/dependency context & Conservative coherence support\\
Application context & Ambiguity explanation and risk reporting\\
\bottomrule
\end{tabularx}
\end{table}

\section{Candidate Generation and Version Selection}

The system first narrows the reference search space and then evaluates
candidate coordinates and versions using multiple evidence families. Library
identity and exact-version confidence are distinct decisions. When adjacent
versions remain inseparable, the system records an ambiguous version set
instead of silently converting a near-tie into an exact fact.

\section{CycloneDX and Vulnerability Design}

Selected components are serialized as CycloneDX components with Package URLs
and evidence properties. Native libraries are represented as file components
with hashes. Vulnerability findings are constrained by component and version
evidence: an uncertain version may support a potential investigation hint, but
not an unqualified confirmed-vulnerability claim.

\section{Experimental Isolation and Fail-Fast Controls}

Alternative reference databases, AI suggestions, R8 diagnostics, and scoring
experiments are isolated from the shared calibrated path. A target improvement
is insufficient for promotion if controlled, clean, partial, no-Maven, or
boundary controls regress.

\section{Use Cases and Behavioral Diagrams}

The following views refine the high-level architecture into actor
responsibilities, runtime behavior, reference construction, deployment
boundaries, and persistent data. Solid blue elements belong to the
authoritative MVP path. Dashed orange elements are optional research or
experiment paths and cannot alter a production decision unless they pass the
promotion gates described previously.

\subsection{Actors and Use Cases}

The principal user is an analyst or auditor who supplies an APK and consumes
the generated evidence. A reference-database maintainer prepares reproducible
Maven examples, while a researcher evaluates changes in isolated databases.
OSV is an external service rather than an internal source of component
identity.

\begin{figure}[H]
\centering
\resizebox{\textwidth}{!}{%
\begin{tikzpicture}[
  actor/.style={
    rectangle,
    rounded corners,
    draw=schoolblue,
    thick,
    fill=blue!5,
    align=center,
    minimum width=31mm,
    minimum height=10mm
  },
  usecase/.style={
    ellipse,
    draw=schoolblue,
    thick,
    fill=white,
    align=center,
    text width=38mm,
    minimum width=45mm,
    minimum height=11mm
  },
  experiment/.style={
    usecase,
    draw=warningorange,
    dashed,
    fill=orange!5
  },
  link/.style={thick,draw=schoolblue},
  explink/.style={thick,dashed,draw=warningorange}
]
\node[actor] (analyst) at (-1.4,4.8) {Analyst / Auditor};
\node[actor] (maintainer) at (-1.4,0.5) {Reference-Database\\Maintainer};
\node[actor] (researcher) at (-1.4,-3.5) {Researcher};
\node[actor] (osv) at (16.0,2.2) {External OSV\\Service};

\node[usecase] (inspect) at (4.6,4.8) {Inspect APK and\\recover evidence};
\node[usecase] (match) at (10.6,4.8) {Match components and\\estimate versions};
\node[usecase] (report) at (10.6,2.2) {Generate CycloneDX\\and reports};
\node[usecase] (builddb) at (4.6,0.5) {Build or update a\\reference database};
\node[usecase] (validate) at (10.2,-0.3) {Validate emitted rows\\against source};
\node[experiment] (experiment) at (4.9,-3.5) {Run isolated target and\\control experiments};
\node[experiment] (promote) at (11.0,-3.5) {Review promotion gates};

\node[
  draw=schoolblue,
  rounded corners,
  inner sep=7mm,
  fit=(inspect)(match)(report)(builddb)(validate)(experiment)(promote),
  label={[schoolblue]above:\textbf{APK-to-SBOM system}}
] (system) {};

\draw[link] (analyst) -- (inspect);
\draw[link] (inspect) -- (match);
\draw[link] (match) -- (report);
\draw[link] (maintainer) -- (builddb);
\draw[link] (builddb) -- (match);
\draw[link] (report) -- (validate);
\draw[link] (report) -- (osv);
\draw[explink] (researcher.east) -- (experiment.west);
\draw[explink] (builddb) -- (experiment);
\draw[explink] (experiment) -- (promote);
\end{tikzpicture}%
}
\caption{Actors and use cases, with experiment-only actions shown as dashed
orange elements.}
\label{fig:pfe-use-cases}
\end{figure}

Figure~\ref{fig:pfe-use-cases} separates ordinary analysis from system
maintenance and research. In particular, building an experimental database is
not equivalent to promoting it: promotion requires an explicit review after
target and control results have been checked.

\subsection{APK-to-Report Activity}

The runtime activity in Figure~\ref{fig:pfe-activity} shows two conservative
branches. A weak or ambiguous match is retained as uncertainty rather than
forced into an exact version, and vulnerability lookup is optional so that
offline runs remain reproducible.

\begin{figure}[H]
\centering
\resizebox{!}{0.68\textheight}{%
\begin{tikzpicture}[
  node distance=5.5mm and 15mm,
  process/.style={
    rectangle,
    rounded corners,
    draw=schoolblue,
    thick,
    fill=blue!5,
    align=center,
    text width=48mm,
    minimum height=9mm
  },
  decision/.style={
    diamond,
    draw=schoolblue,
    thick,
    fill=blue!3,
    aspect=2.3,
    align=center,
    inner sep=1.2mm,
    text width=31mm
  },
  terminal/.style={
    ellipse,
    draw=schoolblue,
    thick,
    fill=evidencegreen!10,
    align=center,
    minimum width=32mm,
    minimum height=8mm
  },
  arrow/.style={-{Latex[length=2mm]},thick,draw=schoolblue},
  branch/.style={font=\scriptsize,fill=white,inner sep=1pt}
]
\node[terminal] (start) {Input APK};
\node[process,below=of start] (inspect) {Inspect archive, manifest, DEX entries,
native files, and forensic indicators};
\node[process,below=of inspect] (analyze) {Parse every recoverable DEX and build
package, class, method, opcode, literal, API, and graph features};
\node[process,below=of analyze] (load) {Load candidate coordinate/version
profiles from the selected SQLite reference database};
\node[process,below=of load] (score) {Compute normalized evidence scores and
rank candidates within each coordinate};
\node[decision,below=8mm of score] (accept) {Identity and version evidence
sufficient?};
\node[process,below left=9mm and 14mm of accept] (uncertain) {Record alternatives,
ambiguity, or abstention; exclude uncertain rows by default};
\node[process,below right=9mm and 14mm of accept] (accepted) {Create accepted
Maven detections with PURLs, confidence, and evidence};
\node[process,below=18mm of accept] (sbom) {Add hash-backed native file
observations and emit CycloneDX 1.5};
\node[decision,below=8mm of sbom] (query) {OSV lookup enabled?};
\node[process,below left=9mm and 14mm of query] (offline) {Continue without
network vulnerability data};
\node[process,below right=9mm and 14mm of query] (osvlookup) {Query OSV for
accepted Maven coordinates and versions};
\node[process,below=18mm of query] (render) {Write matches, SBOM, vulnerability
JSON, and JSON / Markdown / HTML reports};
\node[terminal,below=of render] (end) {Evidence package};

\draw[arrow] (start) -- (inspect);
\draw[arrow] (inspect) -- (analyze);
\draw[arrow] (analyze) -- (load);
\draw[arrow] (load) -- (score);
\draw[arrow] (score) -- (accept);
\draw[arrow] (accept) -| node[branch,near start] {no} (uncertain);
\draw[arrow] (accept) -| node[branch,near start] {yes} (accepted);
\draw[arrow] (uncertain) |- (sbom);
\draw[arrow] (accepted) |- (sbom);
\draw[arrow] (sbom) -- (query);
\draw[arrow] (query) -| node[branch,near start] {no} (offline);
\draw[arrow] (query) -| node[branch,near start] {yes} (osvlookup);
\draw[arrow] (offline) |- (render);
\draw[arrow] (osvlookup) |- (render);
\draw[arrow] (render) -- (end);
\end{tikzpicture}
}
\caption{Activity flow from APK input to the final evidence package.}
\label{fig:pfe-activity}
\end{figure}

\subsection{APK-to-Report Interaction Sequence}

Figure~\ref{fig:pfe-apk-sequence} identifies the runtime ownership of each
operation. The report orchestrator invokes archive inspection and matching;
the matcher itself performs DEX analysis and reads the reference database.
Only accepted decisions are passed to CycloneDX and OSV generation.

\begin{figure}[H]
\centering
\resizebox{\textwidth}{!}{%
\begin{tikzpicture}[
  participant/.style={
    rectangle,
    draw=schoolblue,
    thick,
    fill=blue!5,
    align=center,
    minimum width=27mm,
    minimum height=9mm
  },
  life/.style={draw=gray!65,dashed},
  call/.style={-{Latex[length=1.8mm]},thick,draw=schoolblue},
  ret/.style={-{Latex[length=1.8mm]},thick,dashed,draw=evidencegreen}
]
\node[participant] (user) at (0,0) {Analyst};
\node[participant] (cli) at (3.2,0) {CLI / Report\\Orchestrator};
\node[participant] (extractor) at (6.4,0) {APK Extractor};
\node[participant] (matcher) at (9.6,0) {DEX Analyzer\\and Matcher};
\node[participant] (db) at (12.8,0) {Reference\\SQLite DB};
\node[participant] (output) at (16.0,0) {CycloneDX, OSV,\\Report Writers};

\foreach \p in {user,cli,extractor,matcher,db,output}
  \draw[life] (\p.south) -- ++(0,-82mm);

\draw[call] ($(user.south)+(0,-6mm)$) -- node[above,font=\scriptsize]
  {generate-report(APK, options)} ($(cli.south)+(0,-6mm)$);
\draw[call] ($(cli.south)+(0,-15mm)$) -- node[above,font=\scriptsize]
  {inspect\_apk()} ($(extractor.south)+(0,-15mm)$);
\draw[ret] ($(extractor.south)+(0,-23mm)$) -- node[above,font=\scriptsize]
  {metadata, hashes, native files} ($(cli.south)+(0,-23mm)$);
\draw[call] ($(cli.south)+(0,-32mm)$) -- node[above,font=\scriptsize]
  {match\_apk\_to\_reference\_db()} ($(matcher.south)+(0,-32mm)$);
\draw[call] ($(matcher.south)+(0,-40mm)$) -- node[above,font=\scriptsize]
  {load component feature profiles} ($(db.south)+(0,-40mm)$);
\draw[ret] ($(db.south)+(0,-48mm)$) -- node[above,font=\scriptsize]
  {candidate versions and features} ($(matcher.south)+(0,-48mm)$);
\draw[ret] ($(matcher.south)+(0,-56mm)$) -- node[above,font=\scriptsize]
  {rankings, decisions, ambiguity} ($(cli.south)+(0,-56mm)$);
\draw[call] ($(cli.south)+(0,-64mm)$) -- node[above,font=\scriptsize]
  {accepted detections and APK evidence} ($(output.south)+(0,-64mm)$);
\draw[ret] ($(output.south)+(0,-72mm)$) -- node[above,font=\scriptsize]
  {SBOM and report artifact paths} ($(cli.south)+(0,-72mm)$);
\draw[ret] ($(cli.south)+(0,-80mm)$) -- node[above,font=\scriptsize]
  {evidence package} ($(user.south)+(0,-80mm)$);
\end{tikzpicture}%
}
\caption{Sequence of interactions during an APK-to-report execution.}
\label{fig:pfe-apk-sequence}
\end{figure}

\subsection{Reference-Database Construction Sequence}

The second sequence, Figure~\ref{fig:pfe-reference-sequence}, uses the same DEX
feature representation as APK analysis. This avoids comparing source-level
library descriptions directly with compiled APK bytecode. Conversion records
also preserve skipped and failed artifacts instead of silently treating them
as valid empty libraries.

\begin{figure}[H]
\centering
\resizebox{\textwidth}{!}{%
\begin{tikzpicture}[
  participant/.style={
    rectangle,
    draw=schoolblue,
    thick,
    fill=blue!5,
    align=center,
    minimum width=27mm,
    minimum height=9mm
  },
  life/.style={draw=gray!65,dashed},
  call/.style={-{Latex[length=1.8mm]},thick,draw=schoolblue},
  ret/.style={-{Latex[length=1.8mm]},thick,dashed,draw=evidencegreen}
]
\node[participant] (maintainer) at (0,0) {Database\\Maintainer};
\node[participant] (config) at (3.2,0) {Reference CLI\\and Catalog};
\node[participant] (maven) at (6.4,0) {Maven\\Repository};
\node[participant] (store) at (9.6,0) {Artifact and\\Metadata Store};
\node[participant] (convert) at (12.8,0) {D8 Converter\\and DEX Analyzer};
\node[participant] (sqlite) at (16.0,0) {Selected\\SQLite DB};

\foreach \p in {maintainer,config,maven,store,convert,sqlite}
  \draw[life] (\p.south) -- ++(0,-82mm);

\draw[call] ($(maintainer.south)+(0,-6mm)$) -- node[above,font=\scriptsize]
  {select coordinates and versions} ($(config.south)+(0,-6mm)$);
\draw[call] ($(config.south)+(0,-15mm)$) -- node[above,font=\scriptsize]
  {resolve metadata / download JAR or AAR} ($(maven.south)+(0,-15mm)$);
\draw[ret] ($(maven.south)+(0,-23mm)$) -- node[above,font=\scriptsize]
  {artifact, POM, checksums} ($(config.south)+(0,-23mm)$);
\draw[call] ($(config.south)+(0,-32mm)$) -- node[above,font=\scriptsize]
  {persist bytes and provenance} ($(store.south)+(0,-32mm)$);
\draw[call] ($(store.south)+(0,-40mm)$) -- node[above,font=\scriptsize]
  {convert artifact to comparable DEX} ($(convert.south)+(0,-40mm)$);
\draw[ret] ($(convert.south)+(0,-48mm)$) -- node[above,font=\scriptsize]
  {DEX or explicit skipped / failed record} ($(store.south)+(0,-48mm)$);
\draw[call] ($(store.south)+(0,-56mm)$) -- node[above,font=\scriptsize]
  {extract DEX and POM feature families} ($(convert.south)+(0,-56mm)$);
\draw[call] ($(convert.south)+(0,-64mm)$) -- node[above,font=\scriptsize]
  {transactional component upsert} ($(sqlite.south)+(0,-64mm)$);
\draw[ret] ($(sqlite.south)+(0,-72mm)$) -- node[above,font=\scriptsize]
  {row counts and extraction status} ($(config.south)+(0,-72mm)$);
\draw[ret] ($(config.south)+(0,-80mm)$) -- node[above,font=\scriptsize]
  {auditable build summary} ($(maintainer.south)+(0,-80mm)$);
\end{tikzpicture}%
}
\caption{Sequence used to construct a versioned Maven reference database.}
\label{fig:pfe-reference-sequence}
\end{figure}

\subsection{Component and Deployment View}

The implementation can run through the command-line interface directly or in
the project Docker environment. APKs, selected SQLite databases, and generated
artifacts remain filesystem resources. Network access is required for Maven
acquisition and OSV lookup, but not for deterministic matching with an
existing database. Figure~\ref{fig:pfe-deployment} also makes the authority
boundary explicit: report-only AI and R8 diagnostics may annotate a report but
cannot feed the matcher, CycloneDX component set, or vulnerability labels.

\begin{figure}[H]
\centering
\resizebox{\textwidth}{!}{%
\begin{tikzpicture}[
  component/.style={
    rectangle,
    rounded corners,
    draw=schoolblue,
    thick,
    fill=blue!5,
    align=center,
    minimum width=31mm,
    minimum height=10mm
  },
  storage/.style={
    cylinder,
    shape border rotate=90,
    aspect=0.28,
    draw=schoolblue,
    thick,
    fill=cyan!6,
    align=center,
    minimum width=29mm,
    minimum height=13mm
  },
  external/.style={
    rectangle,
    rounded corners,
    draw=evidencegreen,
    thick,
    fill=green!5,
    align=center,
    minimum width=30mm,
    minimum height=10mm
  },
  experiment/.style={
    rectangle,
    rounded corners,
    draw=warningorange,
    dashed,
    thick,
    fill=orange!5,
    align=center,
    minimum width=34mm,
    minimum height=10mm
  },
  arrow/.style={-{Latex[length=2mm]},thick,draw=schoolblue},
  extarrow/.style={-{Latex[length=2mm]},thick,draw=evidencegreen},
  exparrow/.style={-{Latex[length=2mm]},thick,dashed,draw=warningorange}
]
\node[component] (user) at (0,3.0) {Analyst / Auditor};
\node[component] (cli) at (4.0,3.0) {CLI and Report\\Orchestrator};
\node[component] (extract) at (8.1,4.5) {APK Extractor and\\DEX Analyzer};
\node[component] (matcher) at (8.1,1.5) {Matcher and\\Version Decisions};
\node[component] (writers) at (12.3,3.0) {CycloneDX, OSV, and\\Report Components};

\node[storage] (apkstore) at (4.0,0.0) {APK and Output\\Filesystem};
\node[storage] (shareddb) at (12.3,0.0) {Shared Calibrated\\Reference DB};
\node[external] (mavenrepo) at (16.3,4.5) {Maven\\Repositories};
\node[external] (osvservice) at (16.3,2.2) {OSV API};
\node[experiment] (refbuilder) at (12.3,-3.0) {Reference Builder and\\Isolated DB Experiments};
\node[experiment] (ai) at (5.0,-3.0) {ML/DL, R8, and Version-Band\\Report-Only Diagnostics};

\node[
  draw=schoolblue,
  rounded corners,
  inner sep=5mm,
  fit=(cli)(extract)(matcher)(writers),
  label={[schoolblue]above:\textbf{Local or Docker analysis environment}}
] (runtime) {};

\draw[arrow] (user) -- (cli);
\draw[arrow] (apkstore.north) -- (cli.south);
\draw[arrow] (cli) -- (extract);
\draw[arrow] (extract) -- (matcher);
\draw[arrow] (shareddb.west) -- (matcher.east);
\draw[arrow] (matcher) -- (writers);
\draw[extarrow] (writers) -- (osvservice);
\draw[extarrow] (mavenrepo.east) -- ++(7mm,0) |- (refbuilder.east);
\draw[exparrow] (refbuilder) -- node[right,font=\scriptsize,align=left]
  {promotion only\\after gates} (shareddb);
\draw[exparrow] (ai.west) -- ++(-12mm,0) |- node[near start,left,
  font=\scriptsize,align=left] {annotations\\only} (cli.south);
\end{tikzpicture}%
}
\caption{Component and deployment view with authoritative and report-only
paths separated.}
\label{fig:pfe-deployment}
\end{figure}

\subsection{Reference-Database Data Model}

The physical SQLite schema is centered on one row per Maven coordinate and
version in the \texttt{components} table. Every feature table refers to that
row through \texttt{component\_id}. Figure~\ref{fig:pfe-reference-schema}
groups tables by role to keep the model readable; the labels are the actual
table names used by the implementation.

This one-to-many organization permits transactional replacement of a
component's feature rows while preserving a stable coordinate/version
identity. It also allows new diagnostic feature families to be stored
without changing the meaning of the calibrated production weights. The
schema therefore keeps the authoritative component row stable while letting
feature-specific evidence evolve across extractor revisions, which is exactly
what the thesis needs when it compares calibrated, experimental, and
report-only feature families.

\begin{figure}[H]
\centering
\resizebox{0.98\textwidth}{!}{%
\begin{tikzpicture}[
  entity/.style={
    rectangle,
    draw=schoolblue,
    thick,
    fill=blue!5,
    align=left,
    text width=49mm,
    inner sep=3mm,
    font=\small
  },
  group/.style={
    rectangle,
    rounded corners,
    draw=evidencegreen,
    thick,
    fill=green!4,
    align=left,
    text width=55mm,
    inner sep=3mm,
    font=\small
  },
  diagnostic/.style={
    group,
    draw=warningorange,
    dashed,
    fill=orange!4,
    text width=61mm,
    font=\footnotesize
  },
  relation/.style={-{Latex[length=1.8mm]},thick,draw=schoolblue}
]
\node[entity] (components) at (0,0) {
  \textbf{components}\\
  \texttt{id} (PK)\\
  \texttt{coordinate}, \texttt{group\_id}, \texttt{artifact\_id}\\
  \texttt{version}, \texttt{purl}, \texttt{packaging}\\
  conversion paths, status, and summary counts\\
  \textit{UNIQUE(coordinate, version)}
};

\node[group] (dex) at (-6.7,3.4) {
  \textbf{Artifact inventory}\\
  \texttt{dex\_files}\\
  path, SHA-256, size, class/method counts
};
\node[group] (identity) at (0,4.0) {
  \textbf{Identity features}\\
  \texttt{packages}\\
  \texttt{classes}\\
  \texttt{methods}
};
\node[group] (numeric) at (6.7,3.4) {
  \textbf{Count and shape features}\\
  \texttt{opcodes}\\
  \texttt{opcode\_ngrams}\\
  \texttt{dependency\_shapes}
};
\node[group] (context) at (-6.7,-3.5) {
  \textbf{Relations and context}\\
  \texttt{class\_dependencies}\\
  \texttt{string\_literals}\\
  \texttt{api\_references}
};
\node[group] (pom) at (0,-4.1) {
  \textbf{Declared dependency context}\\
  \texttt{maven\_dependencies}\\
  coordinate, version, scope, optional, type,
  classifier
};
\node[diagnostic] (multi) at (6.7,-3.5) {
  \textbf{Multi-scale diagnostic features}\\
  \texttt{basic\_block\_opcode\_fingerprints}\\
  \texttt{call\_chain\_opcode\_fingerprints}\\
  \texttt{class\_neighbor\_summaries}\\
  \texttt{field\_operation\_sequences}\\
  \texttt{literal\_context\_features}
};

\draw[relation] (components) -- (dex);
\draw[relation] (components) -- (identity);
\draw[relation] (components) -- (numeric);
\draw[relation] (components) -- (context);
\draw[relation] (components) -- (pom);
\draw[relation] (components) -- (multi);
\end{tikzpicture}%
}
\caption{Simplified physical schema of the SQLite reference-feature database;
each arrow represents a \texttt{component\_id} foreign-key relation.}
\label{fig:pfe-reference-schema}
\end{figure}

\section{Conclusion}

The architecture turns APK observations into explicitly bounded SBOM claims.
Its main contribution is not only the final output format, but the discipline
with which each evidence layer is kept separate. The result is a design that
can accept stronger evidence when it exists, preserve uncertainty when it
does not, and still emit a useful CycloneDX-oriented report.

The following chapter explains how these components were implemented and how
their behavior was evaluated. That evaluation matters because the design is
only valuable if the implementation can preserve the same boundaries under
real APKs, controlled demonstrations, and report-only research experiments.
